\section{\vrev{이론적 시사점}}
\label{sec:theoretical-implications}
%% ===================================================================

본 연구의 이론적 기여는 BCM 분야에 \jrev{\chg{i) 관측성 개념의 통합, ii) 불확실성의 확률적 정량화(가우시안과 파레토), iii) 개별 시그모이드 바운딩 이론의 정립, iv) 이중 채널 감쇠 효과의 발견이라는} 네 가지 축}으로 정리된다.
\p{[}표~\vref{tab:ch6-theoretical-contributions}\p{]}는 이론적 기여의 핵심 내용을 종합한다.

\begin{table}[htbp]
\centering
\caption{\jrev{이론적 기여 종합}}
\label{tab:ch6-theoretical-contributions}
\small
\begin{tabularx}{\textwidth}{c l X c}
\toprule
\textbf{\#} & \textbf{기여 영역} & \textbf{핵심 내용} & \textbf{근거} \\
\midrule
1 & 관측성 통합 &
  제어 이론의 관측성 개념을 BCM에 최초 도입;
  $z_O = \kappa \cdot k_O \cdot O_{\mathrm{raw}}$으로 정량화하여
  정적 RTO의 구조적 한계를 극복 &
  H1, \jrev{H3}-a \\
\addlinespace
2 & 시그모이드 바운딩 이론 &
  수정 시그모이드 $S(z) = 2/(1+e^{-z})-1$을 기저 함수로 채택;
  $\DRTO = R_0 + E_{O,bnd} + E_{U,bnd}$로 물리적 경계
  $(0,\; R_{\max})$를 해석적으로 보장 &
  \jrev{H2 ($\kappa = 1.2$ 전제)} \\
\addlinespace
3 & 관측성--불확실성 이원 구조 &
  관측성 효과($E_{O,bnd} \leq 0$)와 불확실성 효과($E_{U,bnd} \geq 0$)의
  구조적 반대 방향 확인; Cohen's $d$: $-1.180$ vs $+1.871$ &
  \jrev{H1, H8} \\
\addlinespace
4 & P85.8 CDF 교차점 규명 &
  가우시안(C2)과 파레토(C3) 분포의 $\eta$ 분위수 함수가
  P85.8에서 교차; 꼬리 효과의 조건부 특성 규명 &
  \jrev{H3}-c \\
\addlinespace
5 & 이중채널 감쇠 발견 &
  P85.8 기준 분할 회귀에서 Tail 영역의 $k_O$ 계수($-0.415$)가
  Core($-0.798$) 대비 $0.52$배 감쇠; 극단 환경에서 관측성
  효과의 상대적 축소 &
  \jrev{H8} \\
\addlinespace
6 & 꼬리 리스크 비선형 증폭 &
  CVaR$_{99}$: C3가 C2 대비 $+24.0\%$;
  분위수 상승에 따른 비선형 격차 확대
  (CVaR$_{90}$: $+20.4\%$ $\to$ CVaR$_{99}$: $+24.0\%$) &
  \jrev{H8} \\
\bottomrule
\end{tabularx}
\end{table}


%% -------------------------------------------------------------------
\subsection{\vrev{\revdel{BCM 이론에 대한 관측성 통합}}}
\label{subsec:observability-integration}
%% -------------------------------------------------------------------

\fdel{전통적 BCM 이론에서 RTO는 BIA(업무영향분석) 시점에 고정되는
정적 매개변수로 취급되어, 시스템 상태의 실시간 변화를 반영하지 못하였다.
본 연구는 제어 이론~\citep{kalman1960control}의 관측성(observability) 개념을
BCM 분야에 최초로 도입하여, 시스템의 내부 상태를 외부 출력으로부터
추정할 수 있는 정도를 $O_{\mathrm{raw}} \in [0,1]$로 정규화하고,
이를 관측성 가중치 $k_O$와 곡률 매개변수 $\kappa = 1.2$를 통해
$\DRTO$의 관측성 조정량으로 변환하였다.}

\fdel{관측성 조정량 $E_{O,bnd} = -R_0 \cdot S(\kappa \cdot k_O \cdot O_{\mathrm{raw}})$는
항상 비양수($\leq 0$)이므로, 관측성이 높을수록 $\DRTO$가 $R_0$ 이하로 단축된다.
이 구조는 \jrev{H3}-a에서 $100\%$ 준수율로 검증되었으며, 이는 확률적 결과가 아닌
모델의 \textbf{구조적 필연}이다. Cohen's $d = -1.180$의 매우 큰 효과 크기는
관측성 통합이 BCM의 복구 예측 정확도를 질적으로 향상시킨다는 이론적 근거를 제공한다.}

\fdel{이러한 관측성 통합은 기존 BCM 이론의 세 가지 공백을 해소한다.
첫째, RTO를 ``고정 상수''에서 ``상태 의존 함수''로 전환함으로써
동적 환경에서의 BCM 의사결정 기반을 확보한다.
둘째, APM(Application Performance Monitoring), SIEM(Security Information and
Event Management) 등 기존 모니터링 인프라의 BCM 활용 경로를 이론적으로 정당화한다.
이는 SRE(Site Reliability Engineering) 분야에서 관측성을 에러 예산(error budget)과
연계하여 서비스 신뢰성을 정량적으로 관리하는 접근(Beyer et al.\ 2016)과
맥락을 공유하며, DevOps 문화에서 모니터링 인프라 투자가 복구 성능을 향상시킨다는
실증적 발견(Forsgren et al.\ 2018; Kim et al.\ 2016)과도 부합한다.
셋째, 관측성 투자($k_O$ 증가)와 불확실성 허용($k_U$ 증가) 간의
trade-off를 $k_O + k_U = 1$ 제약으로 형식화하여,
BCM 자원 배분의 분석적 프레임워크를 제시한다.}


%% -------------------------------------------------------------------
\subsection{\vrev{\revdel{개별 시그모이드 바운딩 이론}}}
\label{subsec:sigmoid-bounding-theory}
%% -------------------------------------------------------------------

\fdel{$\DRTO$ 모델의 수학적 핵심은 수정 시그모이드 함수
$S(z) = 2/(1+e^{-z})-1 \in (-1,1)$에 의한 개별 바운딩(individual bounding)이다.
관측성 효과와 불확실성 효과를 각각 독립적인 시그모이드로 바운딩하여
$\DRTO = R_0 + E_{O,bnd} + E_{U,bnd}$로 정식화함으로써, 다음의 세 가지
이론적 요구사항을 동시에 충족한다.}

% [5단계 Plan H 재편: 구 5.2.2 "개별 시그모이드 바운딩 이론" 3개 항목
%  (물리적 경계 보장 / 무한 미분 가능성 / 교호작용 포착)은 5.2.1 이론적
%  기여로 통합·흡수됨. 기존 \fdel{} 래핑된 본문은 빈 enumerate 번호(1/2/3)
%  잔재를 발생시켜 enumerate 블록 전체 제거.]



%% -------------------------------------------------------------------
\subsection{\vrev{\revdel{이중채널 감쇠 효과와 P85.8 교차점}}}
\label{subsec:dual-channel-crossover}
%% -------------------------------------------------------------------

\fdel{본 연구의 가장 핵심적인 이론적 발견은 P85.8 분할 기준점(CDF 교차점과 일치)을
경계로 하는 이중채널 감쇠(dual-channel attenuation) 효과이다.
C3 조건에서 P85.8 기준으로 분할 회귀를 수행한 결과,
관측성 가중치 $k_O$의 회귀 계수가 Core 영역($\leq$ P85.8)의 $-0.798$에서
Tail 영역($>$ P85.8)의 $-0.415$로 $0.52$배 감쇠된다.}

\fdel{이 발견의 이론적 함의는 세 가지이다.
첫째, 극단적 불확실성 환경에서 관측성 투자의 \textbf{한계효용이 감소}한다.
Core 영역에서 $k_O$를 $0.1$ 증가시키면 $\eta$가 $0.080$ 감소하나,
Tail 영역에서는 동일한 변화가 $\eta$를 $0.042$만 감소시킨다.
이는 BCM 자원 배분에서 극단 시나리오일수록 관측성 투자의
비용-효과성이 감소하며 불확실성 관리가 우선되어야 함을 의미한다.
둘째, 시그모이드 비선형성과 파레토의 \textbf{상호작용 메커니즘}이
규명되었다. Tail 영역에서 큰 $U$ 값이 시그모이드의 포화 구간을 활성화하여
불확실성 채널이 지배적이 되며, $k_O$의 상대적 역할이 축소된다.
셋째, P85.8이 분할 회귀 기준점이자 CDF 교차점으로서 가우시안--파레토 분포 전환의
\textbf{구조적 경계}임이 다중 시드 검증을 통해 확인되어,
이 교차점이 특정 시드에 의존하지 않는 강건한 현상임이 입증되었다.}
